\section{invertibilità locale e funzione implicita}
\begin{theorem}[Lagrange vettoriale]
  \label{th:lagrange}
Sia $\Omega\subset \RR^n$ un aperto convesso, 
sia $\vec f\colon \Omega\to \RR^m$ una funzione differenziabile.
Allora dati $\vec x,\vec y \in \Omega$, $\vec x\neq \vec y$, 
esiste $\vec z \in [\vec x,\vec y]$ (il segmento che congiunge $\vec x$ e $\vec y$)
tale che 
\[
 \abs{\vec f(\vec y)- \vec f(\vec x)}^2
 =  D\vec f(\vec z)[\vec y-\vec x]\cdot (\vec f(\vec y)-\vec f(\vec x)).
\]
In particolare si ha 
\[
\abs{\vec f(\vec y)- \vec f(\vec x)}
\le \max_{\vec z\in [\vec x,\vec y]} \abs{D\vec f(\vec z)} \cdot \abs{\vec y-\vec x}.
\]
\end{theorem}
\begin{proof}
Basta applicare il teorema di Lagrange
alla funzione $g\colon[0,1]\to \RR$
\begin{align*}
  g(t) &= \langle f(\vec x + t(\vec y-\vec x)), f(\vec y)-f(\vec x)\rangle
  g'(t) &= \langle Df(\vec x + t(\vec y - \vec x))[\vec y - \vec x], f(\vec y) - f(\vec x)\rangle
\end{align*}
per ottenere che esiste $t\in [0,1]$ e $z=\vec x + t(\vec y - \vec x)$ 
per cui
\begin{align*}
  \abs{f(\vec y) - f(\vec x)}^2 
    = g(1) - g(0) = g'(t)
    = \langle Df(z)[\vec y - \vec x], f(\vec y) - f(\vec x)\rangle.
\end{align*}
\end{proof}

%%%%%%%%%%%%%%%%%%%%%%%%%%%%%%%%%%%%%%%%%%%%%%%%%%%%%%%%%%%%
\begin{theorem}[invertibilità locale]
%%%%%%%%%%%%%%%%%%%%%%%%%%%%%%%%%%%%%%%%%%%%%%%%%%%%%%%%%%%%
Sia $A$ un aperto di $\RR^n$, $f\in \CC^1(A,\RR^n)$,
$x_0\in A$ e supponiamo che $Df(x_0)$ sia una invertibile.
Allora esiste $\delta>0$ tale che $f(B_\delta(x_0))$ è aperto, 
e la restrizione $f\colon B_\delta(x_0)\to f(B_\delta(x_0))$
è invertibile con inversa di classe $C^1$.
Inoltre per ogni $x \in B_\delta(x_0)$
\[
        D f^{-1}(f(x))=(Df(x))^{-1}.
\]
\end{theorem}
%%%%%%%%%%%%%%%%%%%%%%%%%%%%%%%%%%%%%%%%%%%%%%%%%%%%%%%%%%%%
\begin{proof}
%%%%%%%%%%%%%%%%%%%%%%%%%%%%%%%%%%%%%%%%%%%%%%%%%%%%%%%%%%%%
Sia $M=Df(x_0)^{-1}$. 
Siccome $A$ è aperto e $Df$ è continuo, 
possiamo scegliere $\rho>0$ in modo che
\[    
    \overline{B_\rho(x_0)}\subset A
        \qquad \text{e} \qquad
        \Abs{Df(x)-Df(x_0)} \le \frac{1}{2\Abs M}
                \quad\forall x\in \overline{B_\rho(x_0)}.
\]
Poniamo $r=\frac{\rho}{2\Abs M}$ e
$y_0=f(x_0)$.  
Fissato $\bar y\in B_r(y_0)$
consideriamo l'applicazione
\begin{align*}
        T\colon \overline{B_\rho(x_0)} &\to \RR^n\\
        T(x) &= x + M[\bar y-f(x)].
\end{align*}

Si ha, per $x\in \overline{B_\rho(x_0)}$,
\[
        \Abs{ DT(x)}
        = \Abs{ \mathrm{Id}-M Df(x)} \le 
        \Abs{ M \Abs{\cdot} Df(x_0)-Df(x)}
        \le \frac{1}{2}
\]
e quindi, per il criterio di Lipschitz, la funzione $T$ \`e $\frac 1
2$-Lipschitziana.

Verifichiamo anche che $T(\overline{B_\rho(x_0)})\subset \overline{B_\rho(x_0)}$.
Infatti, per $x\in \overline{B_\rho(x_0)}$, si ha
        \begin{align*}
        |T(x)-x_0|&\le |T(x)-T(x_0)| + |T(x_0)-x_0|\\
        &\le \frac 1 2 |x-x_0| + |M (\bar y-y_0)|
        \le \frac \rho 2 + \Abs{ M}  r
        \le \rho.
        \end{align*}

Dunque $T\colon\overline{B_\rho(x_0)}\to \overline{B_\rho(x_0)}$ è una
contrazione, $\overline{B_\rho(x_0)}$ è uno spazio metrico completo
e quindi, per il Teorema delle contrazioni, 
esiste un unica soluzione dell'equazione $T(x)=x$,
cioé esiste un unico punto $x$ in $\overline{B_\rho(x_0)}$ per cui $f(x)=\bar y$.

Dato $y\in B_r(y_0)$ chiamiamo $g(y)$ l'unico $x\in \overline{B_\rho(x_0)}$
per cui $f(x)=y$. 
Abbiamo quindi definito una funzione $g\colon B_r(y_0) \to \overline{B_\rho(x_0)}$
tale che $f(g(y)) = y$.

Mostriamo ora che $g$ è lipschitziana. Fissati $y,\bar y \in B_r(y_0)$ 
consideriamo la mappa $T(x) = x + Df(x)^{-1}[\bar y - f(x)]$ già considarata prima.
Ricordando che $T$ è $\frac 1 2$-lipschitz, scopriamo che  
\begin{align*}
\frac 1 2 |g(y)-g(\bar y)| 
        &\ge |T(g(y))-T(g(\bar y))| 
        = \abs{g(y) - g(\bar y) + Df(g(\bar y))^{-1}[\bar y-f(g(y))]}\\
        &\ge |g(y)-g(\bar y)| - \Abs{Df(g(\bar y))^{-1}}\cdot \abs{y - \bar y}
\end{align*}
da cui
\[
        \abs{g(y) - g(\bar y)} 
        \le \frac{\abs{y - \bar y}}{2\Abs{ Df(g(\bar y))^{-1}}}
        \le L \abs{y - \bar y}
\]
in quanto $\Abs{Df(x)^{-1}}$ è una funzione continua 
su $\overline{B_\rho(x_0)}$ e dunque ammette minimo (non nullo 
visto che la matrice inversa non è mai nulla).

Dunque per $y \to \bar y$ si ha $g(y) \to g(\bar y)$
e possiamo dunque scrivere la differenziabilità di $f$ in $g(\bar y)$
per $y\to \bar y$, 
\[
   y - \bar y = f(g(y)) - f(g(\bar y)) = Df(g(\bar y))[g(y) - g(\bar y)] + o(\abs{y - \bar y})
\]
da cui si ricava 
\[
  g(y) - g(\bar y) 
  = Df(g(\bar y))^{-1}[y - \bar y + o(\abs{y-\bar y})]
  = Df(g(\bar y))^{-1}[y - \bar y] + o(\abs{y-\bar y})
\]
ovvero $g$ è differenziabile in $\bar y$ e si ha,
come volevamo dimostrare, 
\[
  Dg(\bar y) = (Df(g(\bar y)))^{-1}.
\]
Visto che $Df$ è continua e $g$ è continua anche $Dg$ è continua,
e quindi $g$ è $C^1$.

La dimostrazione è conclusa. 
Per la continuità di $f$ esiste $\delta>0$ per cui $f(B_\delta(x_0)) \subset B_r(y_0)$.
Preso $V=g^{-1}(B_\delta(x_0))$ sappiamo che $V$ è aperto per la continuità
di $g$, e $g\colon V\to B_\delta(x_0)$ è bigettiva e la sua inversa
è la restrizione di $f$ a $B_\delta(x_0)$.
\end{proof}


%%%%%%%%%%%%%%%%%%%%%%%%%%%%%%%%%%%%%%%%%%%%%%%%%%%%%%%%%%%%
\begin{theorem}[funzione implicita, Dini]
%%%%%%%%%%%%%%%%%%%%%%%%%%%%%%%%%%%%%%%%%%%%%%%%%%%%%%%%%%%%
Sia $\Omega$ un aperto di $\RR^n\times \RR^m$ e sia $f\in
\CC^1(\Omega,\RR^m)$. Sia $(x_0,y_0)\in\Omega\subset \RR^n\times \RR^m$.
Se la matrice $D_yf(x_0,y_0)$ delle derivate di $f(x,y)$ rispetto alle
variabili $y$
\[
D_yf(x_0,y_0)=\left(\frac{\partial f^j}{\partial y^k} (x_0,y_0)\right)_{j,k}
\qquad j=1,\ldots,m\quad k=1,\ldots,m
\]
\`e invertibile, allora esiste un intorno $U\subset \RR^n$ di $x_0$
e una funzione $g\in\CC^1(U,\RR^m)$
tale che
\[
        f(x,g(x))=f(x_0,y_0)\qquad \forall x\in U.
\]
Inoltre, per $y=g(x)$ si ha
\[
  Dg(x) = -D_y f(x,y)^{-1} D_y f(x,y).
\]
\end{theorem}
%%%%%%%%%%%%%%%%%%%%%%%%%%%%%%%%%%%%%%%%%%%%%%%%%%%%%%%%%%%%
\begin{proof}
%%%%%%%%%%%%%%%%%%%%%%%%%%%%%%%%%%%%%%%%%%%%%%%%%%%%%%%%%%%%
Consideriamo la funzione $\tilde f\in\CC^1(\Omega,\RR^n\times\RR^m)$
\[
        \tilde f(x,y)=(x,f(x,y)).
\]
si trova facilmente che 
\[      
        D\tilde f(x,y) =
        \left(
        \begin{array}{c|c}
        \mathrm{Id} & 0 \\ \hline
        D_x f & D_y f
        \end{array}
        \right)
\]
dove $\mathrm{Id}$ \`e la matrice identit\`a $n\times n$. 
Dunque essendo $D_y f(x_0,y_0)$ invertibile,
 anche $D\tilde f(x_0,y_0)$ lo \`e e
posso applicare il teorema di invertibilit\`a locale a $\tilde f$.
Dunque esister\`a una funzione $g$ di classe $\CC^1$ in un intorno di
$(x_0,f(x_0,y_0)$ 
tale che $\tilde f(\tilde
g(x,y))=(x,y)$.
Posto $\tilde g(x,y)=(g_1(x,y),g_2(x,y))$ con $g_1$ a valori in
$\RR^n$, e $g_2$ a valori in $\RR^m$ si ha dunque
\[
        (x,y) 
        = \tilde f(g_1(x,y),g_2(x,y))
        = (g_1(x,y),f(g_1(x,y),g_2(x,y)))
\]
da cui si ottiene $g_1(x,y)=x$ e
$f(x,g_2(x,y))=y$.
Baster\`a dunque porre $g(x)=g_2(x,f(x_0,y_0))$ per ottenere $f(x,g(x))=f(x_0,y_0)$.

Per dimostrare la formula di calcolo di $Dg$ \`e sufficiente
differenziare l'equazione
\[
  f(x,g(x))=f(x_0,y_0)
\]
per ottenere
\[
  D_x f(x,g(x)) + D_y f(x,g(x)) Dg(x)=0
\]
da cui la tesi.
\end{proof}

